% !TEX TS-program = pdflatex
% !TeX program = pdflatex
% !TEX encoding = UTF-8
% !TEX spellcheck = fr

\documentclass{beamer}


%\usepackage{fullpage}
%\usepackage[left=2.8cm,right=2.2cm,top=2 cm,bottom=2 cm]{geometry}
\setbeamersize{text margin left=10pt,text margin right=10pt}


\usepackage{amsmath,amssymb} 
\usepackage{fontenc}
\usepackage{textcomp}
\usepackage[utf8]{inputenc}
\usepackage{CJKutf8}
\usepackage[french]{babel}
\usepackage{arabtex,acolor}
\usepackage{txfonts}
\usepackage{tipa}
\usepackage[]{graphicx}
\usepackage{multirow}
\usepackage{hyperref}
\usepackage{colortbl}
\usepackage{tabularray}
\usepackage{listingsutf8}
\usepackage{wrapfig}
\usepackage{multicol}
\usepackage[export]{adjustbox} %for images in table, also for frame
\usepackage[many]{tcolorbox}
\usepackage{wasysym}
\usepackage[french,lined]{algorithm2e}
\usepackage{alltt} %verbatim with commands
\usepackage{longtable}
\usepackage{tabu}


\definecolor{lightblue}{HTML}{D0D2FF}
\definecolor{lightyellow}{HTML}{FFFFAA}
\definecolor{darkblue}{HTML}{0000BB}
\definecolor{olivegreen}{HTML}{006600}
\definecolor{darkgreen}{HTML}{008B45} %009B55
\definecolor{violet}{HTML}{6600CC}
%\definecolor{deeppink}{HTML}{FF1493}
\definecolor{orangey}{HTML}{FFBB00}

\usetheme{Warsaw} % Antibes Boadilla Warsaw

\beamertemplatenavigationsymbolsempty


\let\oldcite\cite
\renewcommand{\cite}[1]{{\bfseries\color{orangey}\oldcite{#1}}}

\tcbuselibrary{listings}

\newcommand{\kurl}[1]{{\scriptsize\bfseries\color{orangey}\url{#1}}}


%\renewcommand{\baselinestretch}{1.5}

\def\supit#1{\raisebox{0.8ex}{\small\it #1}\hspace{0.05em}}

\AtBeginSection{%
	\begin{frame}
		\sectionpage
	\end{frame}
}

\newcommand{\rottext}[2]{%
	\rotatebox{90}{%
	\begin{minipage}{#1}%
		\raggedleft#2%
	\end{minipage}%
	}%
}


\institute{ %
Laboratoire de la Communication dans les Systèmes Informatiques (LCSI)

\vspace{6pt}
École  nationale Supérieure d'Informatique (ESI, ex. INI), Alger, Algérie
}
\author[ \textbf{\footnotesize\insertframenumber/\inserttotalframenumber} \hspace*{1cm} ESI - ARIES Abdelkrime (Master 2022/2023)] %
{ARIES Abdelkrime}
%\titlegraphic{\includegraphics[height=1cm]{../img/esi-logo.png}%\hspace*{4.75cm}~
	
	
	\date{Master 2022/2023} %\today

\titlegraphic{%
	\includegraphics[height=1cm]{../img/esi-logo.png}%
	\hspace{2cm}
	\includegraphics[height=1cm]{../img/lcsi-logo.png}%
}


%\setbeamertemplate{headline}{}



\newcommand{\keyword}[1]{\textcolor{red}{\bfseries\itshape #1}}
\newcommand{\expword}[1]{\textcolor{olivegreen}{#1}}
\newcommand{\optword}[1]{\textcolor{violet}{\bfseries #1}}

\makeatletter
\newcommand\mysphere{%
	\parbox[t]{10pt}{\raisebox{0.2pt}{\beamer@usesphere{item projected}{bigsphere}}}}
\makeatother

%\let\oldtabular\tabular
%\let\endoldtabular\endtabular
%\renewenvironment{tabular}{\rowcolors{2}{white}{lightblue}\oldtabular\rowcolor{blue}}{\endoldtabular}


\NoAutoSpacing %french autospacing after ":"

\def\graphpath{}

\newcommand{\changegraphpath}[1]{\def\graphpath{#1}}


\newcommand{\vgraphpage}[2][.84\textheight]{%
%	\begin{center}%
		\includegraphics[height=#1]{\graphpath #2}%
%	\end{center}%
}

\newcommand{\hgraphpage}[2][\textwidth]{%
%	\begin{center}%
		\includegraphics[width=#1]{\graphpath #2}%
%	\end{center}%
}

\newcommand{\graphpage}[2][]{%
	\includegraphics[#1]{\graphpath #2}%
}

\bibliographystyle{apalike}

\newcommand{\insertbibliography}[2]{
	\appendix
	\section*{Bibliographie}
	\nocite{#2}
%	\makeatletter % to change template
%	\setbeamertemplate{headline}[default] % not mandatory, but I though it was better to set it blank
%	\def\beamer@entrycode{\vspace*{-\headheight}} % here is the part we are interested in :)
%	\makeatother
	\begin{multicols*}{2}[\frametitle{\insertsection} \usebeamertemplate*{frametitle}]%\usebeamertemplate*{frametitle}\frametitle{Références}
		\tiny
		\bibliography{#1}
	\end{multicols*}
}

\definecolor{my-grey}{RGB}{233, 233, 233}

\newcommand{\insertlicence}{
	\begin{frame}[plain]
	\frametitle{Licence : CC-BY 4.0}
%	\framesubtitle{Licence: CC-BY-NC 4.0}

	\begin{tcolorbox}[colback=cyan,
		colframe=cyan,  
		arc=0pt,outer arc=0pt,
		valign=top, 
		halign=center,
		width=\textwidth]
		
		\includegraphics[width=.5cm]{../img/licence/cc_icon_white_x2.png}
		\includegraphics[width=.5cm]{../img/licence/attribution_icon_white_x2.png}
		
		\color{white}
		\bfseries Attribution 4.0 International (CC BY 4.0) \\
		\tiny \url{https://creativecommons.org/licenses/by/4.0/deed.fr}
		
	\end{tcolorbox}\vspace{-.5cm}
	\begin{tcolorbox}[colback=my-grey,
		colframe=my-grey,  
		center, arc=0pt,outer arc=0pt,
		valign=top, 
		halign=left,
		width=\textwidth]
		
		\tiny
		
		\begin{center}
			\bfseries\Large
			Vous êtes autorisé à :
		\end{center}
		
		\begin{minipage}{0.83\textwidth}
			\begin{itemize}
				\item[] \textbf{Partager} — copier, distribuer et communiquer le matériel par tous moyens et sous tous formats
				\item[] \textbf{Adapter} — remixer, transformer et créer à partir du matériel
				pour toute utilisation, y compris commerciale.
			\end{itemize}
		\end{minipage}
		\begin{minipage}{0.15\textwidth}
			\includegraphics[width=\textwidth]{../img/licence/FreeCulturalWorks_seal_x2.jpg}
		\end{minipage}
	
		
		\begin{center}
			\bfseries\Large
			Selon les conditions suivantes :
		\end{center}
		
		\begin{itemize}
			\item[] \textbf{Attribution} — Vous devez créditer l'Œuvre, intégrer un lien vers la licence et indiquer si des modifications ont été effectuées à l'Oeuvre. Vous devez indiquer ces informations par tous les moyens raisonnables, sans toutefois suggérer que l'Offrant vous soutient ou soutient la façon dont vous avez utilisé son Oeuvre. 
			\item[] \textbf{Pas de restrictions complémentaires} — Vous n'êtes pas autorisé à appliquer des conditions légales ou des mesures techniques qui restreindraient légalement autrui à utiliser l'Oeuvre dans les conditions décrites par la licence.
		\end{itemize}
		
	\end{tcolorbox}
	
%	\begin{center}
%		\bfseries Attribution 4.0 International (CC BY 4.0)
%		\url{https://creativecommons.org/licenses/by/4.0/deed.fr}
%	\end{center}

%	\tiny
%
%	Vous êtes autorisé à : 
%	\begin{itemize}
%		\item \textbf{Partager} — copier, distribuer et communiquer le matériel par tous moyens et sous tous formats
%		\item \textbf{Adapter} — remixer, transformer et créer à partir du matériel
%	\end{itemize}
%	
%	Selon les conditions suivantes : 
%	\begin{itemize}
%		\item \textbf{Attribution} — Vous devez créditer l'Œuvre, intégrer un lien vers la licence et indiquer si des modifications ont été effectuées à l'Oeuvre. Vous devez indiquer ces informations par tous les moyens raisonnables, sans toutefois suggérer que l'Offrant vous soutient ou soutient la façon dont vous avez utilisé son Oeuvre.
%		\item \textbf{Pas d'Utilisation Commerciale} — Vous n'êtes pas autorisé à faire un usage commercial de cette Oeuvre, tout ou partie du matériel la composant. 
%		\item \textbf{Pas de restrictions complémentaires} — Vous n'êtes pas autorisé à appliquer des conditions légales ou des mesures techniques qui restreindraient légalement autrui à utiliser l'Oeuvre dans les conditions décrites par la licence.
%	\end{itemize}

	\end{frame}
}

\settowidth{\leftmargini}{\usebeamertemplate{itemize item}}
\addtolength{\leftmargini}{\labelsep}

\AtBeginDocument{
	\newcolumntype{L}[2]{>{\vbox to #2\bgroup\vfill\flushleft}p{#1}<{\egroup}} 
	
	\begin{frame}[plain]
		\maketitle
	\end{frame}

	\insertlicence
}


% needs etoolbox; to break links after -
\appto\UrlBreaks{\do\-}


\makeatletter
\newcommand{\xRightarrow}[2][]{\ext@arrow 0359\Rightarrowfill@{#1}{#2}}
\makeatother


\usefonttheme{structurebold}
%\usefonttheme{professionalfonts}

\title[TALN : Présentation] %
{Traitement automatique du langage naturel\\Présentation du module} 

\begin{document}


\begin{frame}
\frametitle{Pré-requis}

\begin{itemize}
	\item Théorie des langages et programmation (THP)
	\item Probabilités et statistiques
	\item Algèbre et analyse 
	\item Apprentissage automatique 
	\item Recherche opérationnelle
	\item Connaissance basique sur les langues
\end{itemize}

\end{frame}


\begin{frame}
\frametitle{Objectifs}

\begin{itemize}
	\item Apprendre quelques notions linguistiques ;
	\item Appliquer des concepts vus dans les modules prérequis afin de résoudre des taches TALN ;
	\item Découvrir quelques outils et ressources du TALN ;
	\item Savoir les axes de recherche dans ce domaine ;
	\item Apprendre comment décrire une nouvelle méthode TALN et la tester dans le contexte d'un article scientifique.
\end{itemize}

\end{frame}


\begin{frame}
\frametitle{Informations}

\begin{itemize}
	\item 5 jours (24h)
	\item Cours (15h)
	\begin{itemize}
		\item Cours interactifs (questions pendant le cours)
		\item QCM à la fin de chacun des 3 premiers jours
	\end{itemize}
	\item TP (9h)
	\begin{itemize}
		\item 3 TPs notés
		\item Découverte de quelques outils TALN
		\item Réalisation d'une application TALN
		\item Rédaction d'un article sur cette application
	\end{itemize}
	\item Examen écrit 
	\item Note = 20\% (QCM) + 30\% (TP) + 50\% (Examen)
\end{itemize}

\end{frame}

\begin{frame}
\frametitle{Programme}

\scriptsize
\begin{tblr}{
		colspec = {p{.12\textwidth}lp{.37\textwidth}lp{.38\textwidth}},
		row{odd} = {lightblue},
		row{1} = {darkblue, fg=white, font=\bfseries, valign=m, halign=c},
		column{2,4}={white},
		column{1}={bg=darkblue, fg=white, font=\bfseries, valign=m},
		row{even} = {white},
		cell{1}{1}={white},
		colsep=3pt,
		rowsep=0pt,
	}


%	\hline
	&& 9H00 - 12H00 && 13H00 - 16H00 \\
	
	&&&&\\
	
%	\hline
	Dimanche \newline 20/11/2022 && \textbf{CM01.} Introduction au TALN \newline 
	\textbf{CM02.} Traitements basiques du texte &&  / \\
	
	&&&&\\
	
%	\hline
	Lundi \newline 21/11/2022 && \textbf{CM03.} Apprentissage automatique pour TALN &&  
	\textbf{CM04.} Modèles de langage \newline \textbf{CM05.} Étiquetage morpho-syntaxique \\
	
	&&&&\\
	
%	\hline
	Mardi \newline 22/11/2022 && \textbf{CM06.} Sémantique lexicale &&  / \\
	
	&&&&\\
	
%	\hline
	Mercredi \newline 23/11/2022 && \textbf{TP01.} Découverte de quelques outils TALN &&  
	\textbf{TP02.} Implémentation d'une application TALN \\
	
	&&&&\\
	
%	\hline
	Jeudi \newline 24/11/2022 && \textbf{CM07.} Quelques applications &&  
	\textbf{TP03.} Rédaction d'un article concernant cette application \\
	
%	&&\\
%	
%	\hline
\end{tblr}

\end{frame}

\begin{frame}
	\frametitle{Préparation}
	
	\begin{itemize}
		\item Environnement de développement (au choix)
		\begin{itemize}
			\item Google Colab (il faut une connexion Internet)
			\item Sur machine :
			\begin{itemize}
				\item Installer Anaconda ({\tiny \url{https://www.anaconda.com/products/distribution\#Downloads}})
				\item Installer NLTK (Anaconda navigator)
				\item Installer tensorflow (Anaconda navigator)
			\end{itemize}
		\end{itemize}
		\item Environnement \LaTeX\ pour l'édition de l'article (au choix)
		\begin{itemize}
			\item Overleaf (il faut une connexion Internet)
			\item Sur machine :
			\begin{itemize}
				\item Installer \LaTeX ({\tiny \url{https://www.latex-project.org/get/}})
				\item Installer TexStudio ({\tiny \url{https://www.texstudio.org}})
			\end{itemize}
		\end{itemize}
	\end{itemize}
	
\end{frame}

\begin{frame}
	\frametitle{Cours similaires}
	
	\begin{itemize}
		
		\item Natural Language Processing S21 (à jours).
		Carnegie Mellon University. 
		\url{http://demo.clab.cs.cmu.edu/NLP/}
		
		\item CS388: Natural Language Processing (2018). 
		University of Texas. 
		\url{https://www.cs.utexas.edu/~mooney/cs388/}
		
		\item CSE 517: Natural Language Processing ().
		University of Washington.
		\url{https://courses.cs.washington.edu/courses/cse517/}
		
		\item CS 294-5: Statistical Natural Language Processing (2005). 
		University of Berkely. 
		\url{https://people.eecs.berkeley.edu/~klein/cs294-5/index.html}
		
		\item CS224n: Natural Language Processing with Deep Learning (à jours).
		Stanford university.
		\url{https://web.stanford.edu/class/cs224n/}
		
		\item CS 5340/6340: Natural Language Processing (à jours). 
		University of Utah.
		\url{https://my.eng.utah.edu/~cs5340/schedule.html}
		
	\end{itemize}
	
\end{frame}

\begin{frame}
	\frametitle{A lire (NLP)}
	
	\begin{itemize}
		
		\item \textbf{Traitement automatique du langage naturel - support de cours}
		
		par Abdelkrime Aries
		
		\url{https://github.com/projeduc/ESI_2CS_TALN/releases/download/4.0/taln_v2.pdf}
		
		
		\item \textbf{Speech and Language Processing} 
		
		par Dan Jurafsky et James H. Martin
		
		\url{https://web.stanford.edu/~jurafsky/slp3/}
		
		\item \textbf{Natural langguage processing}
		
		par Jacob Esenstein 
		
		\url{https://cseweb.ucsd.edu/~nnakashole/teaching/eisenstein-nov18.pdf}
		
		\item \textbf{Modern Approaches in Natural Language Processing - support de cours (LMU Munich)}
		
		par Carolin Becker, et al.
		
		\url{https://slds-lmu.github.io/seminar_nlp_ss20/index.html}
		
		
		
	\end{itemize}
	
\end{frame}

\begin{frame}
	\frametitle{A lire (ML)}
	
	\begin{itemize}
		
		
		\item \textbf{Dive into Deep Learning} (online)
		
		\url{https://d2l.ai/index.html}
		
		\item \textbf{DEEP LEARNING, support de cours NYU CENTER FOR DATA SCIENCE} 
		
		par Yann LeCun et  Alfredo Canziani
		
		\url{https://atcold.github.io/pytorch-Deep-Learning/}
		
		
	\end{itemize}
	
\end{frame}

\begin{frame}
	\frametitle{Clause de non-responsabilité}
	
	\begin{tcolorbox}[colback=red,
		colframe=red,  
		arc=0pt,outer arc=0pt,
		valign=top, 
		halign=center,
		width=\textwidth]
		
		\color{white}
		{\bfseries Clause de non-responsabilité}
		
	\end{tcolorbox}\vspace{-.5cm}
	\begin{tcolorbox}[colback=my-grey,
		colframe=my-grey,  
		center, arc=0pt,outer arc=0pt,
		valign=top, 
		halign=left,
		width=\textwidth]
		
		\begin{itemize}
			\item Un ingénieur doit être une force du bien ; il ne faut pas utiliser notre connaissance pour nuire aux autres.
			\item Ce cours n'encourage pas l'utilisation malicieuse des techniques TALN comme l'espionnage, la collecte illégale d'information, la mis-information via des chatbots, etc.
			\item Tout exemple qui peut contenir une de ces applications est purement éducatif.	
		\end{itemize}
		
		
	\end{tcolorbox}
	
\end{frame}


\end{document}

